\SourceLeaf{192}liche Auflösung des Gordischen Knotens. Inzwischen kann man zugeben, dass diese Auslassung
hier in so fern wenig auf sich hat, als die Beweisart des Verf. für die Unstatthaftigkeit des
dritten Falls, wenn sie zulässig wäre, auch auf diesen Fall von selbst erstreckt werden könnte.
Allein eben diesem angeblichen Beweise der Unstatthaftigkeit des dritten Falls können wir keine
Gültigkeit zugestehen. Der Verf. stellt die Sache so vor. Wenn z. B. in dem dritten Falle
angenommen wird, die beiden ersten Stücke seien gleich, das dritte aber grösser, so wäre
$DC$ also grösser als $CL$. Da nun aber $AML$ gleichfalls ein gleichschenkliges Dreieck ist,
dem dieselbe Grundbedingung zukommt, wie dem ursprünglichen Dreieck $AON$, so müsste, wenn
jener dritte Fall mit seiner angenommenen Unterabtheilung der gültige wäre, $DC=CL$ sein,
in Widerspruch mit dem vorher Gefundenen.

Wir haben, wie wir glauben, bei diesem Moment des Beweises das, worauf es ankommt, noch etwas
klarer und bestimmter nach der Ansicht des Verf. angedeutet, als er es selbst gethan hat,
wodurch dann aber auch die Schwäche desselben, wie uns scheint, leichter erkannt wird. Denn
offenbar ist hier ganz willkürlich angenommen, dass bei allen gleichschenkligen Dreiecken mit
dem Winkel $A$ an der Spitze und grösserm Winkel an der Basis, wenn mit ihnen die im Lehrsatz
angezeigte Construction vorgenommen wird, die Folge der abgeschnittenen Stücke in Rücksicht
auf ihr Gleichbleiben, Grösser- oder Kleinerwerden, allemal, unabhängig von der Grösse der
Seiten, nothwendig dieselbe sein müsse, eine Annahme, die doch unmöglich als von selbst evident
betrachtet werden darf. Da sich nun aber hierauf allein der versuchte Beweis der
Unstatthaftigkeit des dritten (wie auch vierten und fünften) Falls stützt, und der ganze Artikel
auch keine andere Ressourcen zum Beweise der Unstatthaftigkeit des übergangenen Falls
darbietet, so glauben wir hierdurch das oben ausgesprochene Urtheil hinlänglich gerechtfertigt
zu haben, wobei wir aber gern der ganzen übrigen sinnreichen Durchführung in den folgenden
Artikeln volle Gerechtigkeit widerfahren lassen.

\begin{center}\rule{3cm}{0.4pt}\end{center}

\Needspace{10\baselineskip}
\begin{center}\textsc{Bemerkung.}\end{center}

{\small
Diese bereits in Band IV, Seite 368 bis 370 abgedruckte Anzeige ist hier der Vollständigkeit
wegen reproducirt worden. Die Figur, die sich weder beim Originale noch bei dem ersten
Abdrucke findet, soll das Verständniss des Textes erleichtern.

\begin{flushright}\textsc{Stäckel.}\end{flushright}
}

\begin{center}\rule{3cm}{0.4pt}\end{center}
